% Ad M. Brutum 1.9 (ad-brutum-01-09)
% Source: https://raw.githubusercontent.com/PerseusDL/canonical-latinLit/master/data/phi0474/phi059/phi0474.phi059.perseus-lat1.xml
% Fetched by scripts/fetch_latin.py

% Dateline (Perseus): Scr. Romae. in. Quint. ...ut videtur a. 711 (43) .

\ciceroLetterOpener{Cicero Bruto salutem}

\ciceroSection{1} fungerer eo officio quo tu functus es in meo luctu teque per litteras consolarer, nisi scirem iis remediis quibus meum dolorem tu levasses te in tuo non egere, ac velim facilius quam tunc mihi nunc tibi tute medeare. est autem alienum tanto viro quantus es tu, quod alteri praeceperit id ipsum facere non posse. me quidem cum rationes quas conlegeras tum auctoritas tua a nimio maerore deterruit. Cum enim mollius tibi ferre viderer quam deceret virum praesertim eum qui alios consolari soleret, accusasti me per litteras gravioribus verbis quam tua consuetudo ferebat.

\ciceroSection{2} itaque iudicium tuum magni aestimans idque veritus me ipse conlegi et ea quae didiceram, legeram, acceperam, graviora duxi tua auctoritate addita. ac mihi tum, Brute, officio solum erat et naturae, tibi nunc populo et scaenae, ut dicitur, serviendum est. nam cum in te non solum exercitus tui sed omnium civium ac paene gentium coniecti oculi sint, minime decet propter quem fortiores ceteri sumus eum ipsum animo debilitatum videri. quam ob rem accepisti tu quidem dolorem (id enim amisisti cui simile in terris nihil fuit), et est dolendum in tam gravi vulnere ne id ipsum, carere omni sensu doloris, sit miserius quam dolere, sed ut modice ceteris utile est, ita tibi necesse est.

\ciceroSection{3} scriberem plura nisi ad te haec ipsa nimis multa essent. nos te tuumque exercitum exspectamus; sine quo, ut reliqua ex sententia succedant, vix satis liberi videmur fore. de tota re publica plura scribam et fortasse iam certiora iis litteris quas veteri nostro cogitabam dare.
